% Conclusion, limitations and deviations. Every number is a macro; process facts from process.tex.
% Deviations D1-D7: plan.json design.deviations.
\section{Conclusion}
\label{sec:conclusion}

Two years after California's \$20 fast-food minimum wage took effect, a state-level QCEW comparison
whose specification was chosen by a pre-specified rule, with inference built for one treated
state, does not determine the sign of its effect on limited-service restaurant employment. The
primary estimate is $\HTwoPct$ percent, with a 95 percent interval from $\HTwoPctLo$ to
$\HTwoPctHi$ percent; the re-estimated specifications of both published camps are also
inconclusive; and employment outside limited-service restaurants deviated from its counterfactual
by a similar amount in the opposite direction, so the estimate is not separable from a
California-wide shock. The first stage in weekly earnings is not supported at the pre-specified
critical value, and two post hoc checks of that verdict pull in opposite directions; the change
between the two years is uninformative.

The main lesson is about the cost of conservative inference with one treated state, and it was
known before the post-period data were read. With the larger of the scaled and unscaled placebo
standard errors, chosen to protect the level of California's tests, and the few effective degrees
of freedom of a placebo standard error, the 90 percent interval has a half-width of
$\HNinetyHTwo$ log points, and this design cannot tell a decline of about 3.5 percent from no
effect: ``inconclusive'' had probability $\ProbTauAltUndUndet$ under the first and
$\ProbTauZeroUndUndet$ under the second. The published estimates that disagree are both consistent
with our interval. The maximum is far from tight for California's employment statistics, and with
the scaled standard error alone the interval would lie below zero. That standard error over-rejects
if shocks have a state-level component, the threat against which the maximum was chosen a priori;
the pre-2020 diagnostics are mixed on whether such a component is present (one population tercile
supports it; the decision set, California's own tests and the Ferman--Pinto fit do not), and the
paper keeps the pre-specified maximum. Separating the published readings would need inference that is valid for California
at a smaller cost in width, more treated units, or more precise data.

\section*{Limitations and deviations from the approved design}
\label{sec:limitations}

\paragraph{Limitations.}
\begin{enumerate}
  \item \emph{Precision.} The design could not reach ``smaller than 2\%'' (Section~\ref{sec:res-power}):
        the pre-period half-width $h_{90} = \HNinetyHTwo$ exceeds $b = 0.02$, and the probability
        of ``inconclusive'' was $\ProbTauZeroUndUndet$ with no effect. The equivalence test therefore
        cannot separate the published readings, which was part of the study's motivation. The width
        is set mainly by the pre-specified maximum standard error, which overstates California's
        pre-period spread ($k = \KDiagHTwo$); with the scaled standard error alone, which
        over-rejects if shocks have a state-level component, the H2 interval would lie below zero.
        The evidence on whether that component is present is mixed
        (Sections~\ref{sec:res-inference} and~\ref{sec:res-h2}).
  \item \emph{The first stage.} H1's critical value comes from the decision-set quantile, not from
        $t_\nu$; the post hoc checks (Section~\ref{sec:res-posthoc}) show that it is set by one state
        and that the selected estimator's pre-2020 wage counterfactual drifts, and they cannot be
        resolved within this design. The selected estimator was chosen on employment and fits
        California's pre-period wages poorly. Weekly earnings per job are not hourly wages, and wages
        are nominal.
  \item \emph{Identification.} The placebo outcome fails, so a California-wide deviation cannot be
        separated from the effect. The selection rule favours candidates that track California's own
        2014--2022 minimum-wage ramp, and its winner's employment pseudo-ATTs are all positive; the
        out-of-fit 2014--2019 employment drift (post hoc) points in the same direction as the placebo
        failure. Neither is corrected for.
  \item \emph{Inference.} California's level is argued from the conservative standard error, not
        measured: no placebo state is as large. If California's own noise exceeds what its size
        implies, the tests over-reject; in the Monte Carlo design with California twice as noisy as
        its donors, the selected candidate's two-sided test at $t_J$ rejected in a share
        $\McDDTPOPWERFiftyHeteroACATwoRhoZeroZeroMaxTJTwo$ of replications. The power quantities rest
        on 2014--2019 noise; the post-period recomputations are reported next to them.
  \item \emph{Measurement.} The estimand is an intent-to-treat effect on QCEW 722513 jobs, covered
        and uncovered establishments together; covered chains are also coded in 722515 and 722511,
        and QCEW 722513 may track covered-chain employment poorly across states and over time
        \citep{melo2026covered}. The QCEW counts jobs, not hours. 2026Q1 is preliminary; the
        final-data window drops it.
  \item \emph{Unexplained cross-check differences.} Our fixed-horizon replication of Clemens and
        coauthors differs from their print by $\CemnFhSimpleDiff$ and $\CemnFhDetrDiff$ index points,
        and our SDID under Dube's timing differs from his by $\DubeTimingEmpDiff$ ($100\times$ log
        points); neither difference was traced.
  \item \emph{Not blind.} The PI and the literature scout knew the published estimates and a public
        descriptive series through March 2026 before the design was fixed
        (Section~\ref{sec:prespec}). Several rules that settle the outcomes of the code check were
        adopted after the last of the \CrewRoundsDesign{} design review rounds.
\end{enumerate}

\paragraph{Deviations from the approved design.} The first four are departures the approved design
itself makes from the track rules and the published specifications (Section~\ref{sec:departures});
the last four arose while the runs were made or the paper was written.
\begin{enumerate}
  \item[D1.] Standard errors are not clustered at the level of treatment: with one treated state,
        cluster-robust standard errors are invalid; placebo-based inference with a pre-specified
        level check is used instead.
  \item[D2.] S1 adapts the monthly fixed-horizon estimator of Clemens and coauthors to a pooled
        quarterly average, which needs seasonal adjustment; the fixed-horizon number is reported as a
        cross-check. S1 keeps their endpoint trend, whose start, 2021Q3, is the only COVID quarter
        used in any estimate.
  \item[D3.] S2 uses our implementation of Algorithm~1 of \citet{arkhangelsky2021synthetic}, with an
        intercept and separate penalties, and our seasonal factors, not Dube's code; the cross-check
        measures the combined difference.
  \item[D4.] R50 keeps Alaska and Michigan, which Clemens and coauthors drop; the cross-check against
        their print uses their 48-state pool.
  \item[D5.] The design gave no yearly H1 statistic. Recorded after the first code review and before
        any analysis run, the yearly wage tests use the pooled H1 critical values and $\nu$, so they
        are approximate. They decide whether yearly elasticities are reported; none is.
  \item[D6.] The frozen requirements pin the released \texttt{wrapt} 2.5.2 rather than the
        pre-release the design named, on the PI's instruction; no estimate depends on it.
  \item[D7.] Post hoc and report-only: run \cand{posthoc} (Section~\ref{sec:res-posthoc} and the
        employment drift in Section~\ref{sec:rob-event}) was added after the results review had
        seen every post-period result. Its code passed the last of the \CrewRoundsCode{} code review
        rounds; that round's minor findings were then fixed and, with the human PI's agreement, the
        run was made without a further review, so no code review covered the code it ran; the second
        results-review round read that code, the same file that ran, after the run. The ledger counts it among
        the analysis runs not covered by a passing code review ($\CrewRunsUnreviewed$ in all). It reports facts
        about Georgia's 2014--2015 series, the H1 critical values without Georgia's decision-set tests
        and with $\nu$ recomputed, a leave-one-out over the ten decision-set states, and the
        2014--2019 drift of the event path. It changes no label. Known defect: its leave-one-out
        change flags (\texttt{b\_loo\_*\_changes\_chi}, \texttt{b\_loo\_n\_changing},
        \texttt{b\_loo\_changing\_states} and the matching line of its log) all read 1, because they
        compare values stored at ten significant digits with a tolerance of $10^{-12}$; no claim uses
        them, the statement that only Georgia moves the critical value rests on the critical values
        themselves, and a clean rerun reproduces the same flags.
  \item[D8.] Descriptive and report-only: run \cand{describe}, added at the paper stage after the
        paper review asked for a summary-statistics table, produces only the levels of
        Table~\ref{tab:sumstats} (means of employment, average weekly wage and establishments for
        California and the two pools, and sample sizes) and no statistic comparing California with
        another unit or one period with another. The code review rounds had been used, so it was
        run without a code review, with the human PI's agreement; the ledger counts it among the
        analysis runs not covered by a passing code review. It feeds the table only and changes no
        label.
\end{enumerate}

\paragraph{How the results were checked.} The results were reviewed in \CrewRoundsResults{} rounds
by a reviewer agent that did not produce them, before the human PI approved them. The ledger holds
\CrewRuns{} runs (\CrewRunsAnalysis{} analysis runs); \CrewRunsSuperseded{} were replaced by a later
run with the same identifier and are archived. Clean rerun of every run from scratch, as recorded
when this version was compiled: \CrewRerunStatus{}, with \CrewRerunReproduced{} of
\CrewRerunTotal{} runs reproduced.
